\section{The gaps of the Eulerian polynomial}
Let $A_m(t)=\sum_{\sigma\in S_m}t^{\des(\sigma)}$.
These polynomials are monic, with $A_m(0)=1$ and $m-1$ simple negative zeros; this follows by induction from
\[
 A_m=(1+(m-1)t)A_{m-1}+t(1-t)A'_{m-1}.
\]
The earlier common intervals of Theorem~\ref{main} can now be chosen without reference to a dyadic or trigonometric mesh.

\begin{theorem}\label{gapthm}
Let $n\ge6$ and let $\xi_1<\cdots<\xi_{n-2}<0$ be the zeros of $A_{n-1}$.
For every $3\le k\le n$, each interval $(\xi_i,\xi_{i+1})$ contains precisely one zero of $D_{n,k}(t)/t$.
These $n-3$ zeros are simple and exhaust its zeros with multiplicity.
The same statement holds for $\sum_{k=3}^n w_kD_{n,k}(t)/t$ whenever all $w_k\ge0$ and some $w_k>0$.
\end{theorem}
The simple zero of the unmodified row at the origin is accounted for in Theorem~\ref{main}. The statement about mixtures involving $k=1,2$ is Theorem~\ref{allcuts}; those cutoffs are outside the Eulerian-gap assertion. The double zero in rank five prevents extending this assertion to all smaller ranks.

We first give the extra estimate needed to work at an Eulerian zero, where the predecessor's leading term cancels.
Retain the notation $f_n(q)=\sum_{h=1}^n a_hq^h$, $a_n=p_n$ and $E(v)$ from above.

\begin{lemma}\label{invstrong}
For $n\ge128$ and $3\le k\le n$, $a_{n-1}=0$ and $E(4/5)<1$. The same estimates hold for the reversed row.
\end{lemma}
\begin{proof}
Use the positive-loss logarithm $B$ of the alphabet product defined above. A term with length $jb$ has loss $jb-j/d$. Every nonprincipal term loses at least half its length, and the smallest possible loss is two. The assertions about degree $n-1$ and reversal follow. We give explicit upper bounds for the exponential of the logarithm; they also specify finite rational checks for all constants. Write $v=4/5$, $V'=93/100$ and $s=179/200$.

For $k\ge58$, the total mass of the logarithm from a length $j\ge k$ at radius $v$ is at most $(21/40)s^j$. Indeed the primitive terms cost at most $s^j/2$ by counting the divisors $d>1$ and using $v\le s^2$. The repetitions cost at most $v^j/[2(1-v^j)]\le v^j\le s^j/40$. Thus the total mass is at most $c=5s^k$. At $k=58$ one has $(2k+1)c<1$, and this expression decreases in $k$. The nonempty terms of the exponential then cost at most $e^c-1\le c/(1-c)<1/(2k)\le p_n$, even if each residual permutation fraction is replaced by one.

For $8\le k\le57$, let $C(w)$ bound the mass of the logarithm for all $j\ge8$. For $8\le j<200$ sum
\[
 \frac1j\sum_{d\in I_j,\ d>1} w^{j-j/d}
 +\frac1{2j(1-w^8)}\sum_{d\in I_j} w^{2j-j/d},
 \quad I_j=\{d:d\mid j,\ 4\nmid d,\ 9\nmid d\}.
\]
This includes every square-free divisor. For the remaining $j$ append $3u^{200}/[2(1-u)]$, with $u=179/200$ at $w=v$ and $u=193/200$ at $w=V'$. This bounds the primitive and repeated terms as above using $w\le u^2$. The sums give $C(v)<0.365$ and $C(V')<2.55$. A positive residual length contributes $p_{n-i}\le2p_n$. A term filling the rank has loss at least $\lceil n/2\rceil$ and $1/p_n\le2k$. Consequently
\[
 E(v)\le 2(e^{0.365}-1)+114(v/V')^{64}(e^{2.55}-1)<0.971.
\]

For $3\le k\le7$ first collect the signed logarithm at length $i$ and alphabet degree $h$. For $h\mid i$, $h<i$, its coefficient is
$i^{-1}\sum_{d\mid i/h,\ hd\ge k}\mu(d)$: write $j=hd$ and $b=i/(hd)$ in $B$. It vanishes for $h\ge k$. Define $C_k(w)$ by summing the absolute collected coefficients times $w^{i-h}$ for $i<200$ and adding
$(k-1)^2w^{200-(k-1)}/[200(1-w)]$. This tail uses $h<k$ and replaces the divisor sum by minus its complementary sum over $hd<k$. Direct evaluation gives $C_k(v)<0.672$ and $C_k(V')<1.573$ throughout these five cutoffs. The recurrence for the fractions gives $m p_m=(m-1)p_{m-1}+p_{m-k}$, for $m\ge k$. The maximum-to-minimum ratio on the $k$ entries starting at $m=20$ is at most $201/200$ for each cutoff. The recurrence propagates this ratio to all $m\ge20$. Thus for residual length at least 20, relative to $p_n$, the weight is at most $201/200$. Otherwise the used length is at least 109, the loss at least 55, and the relative weight at most $2k$. This yields
\[
 E(v)\le(201/200)(e^{0.672}-1)+14(v/V')^{55}(e^{1.573}-1)<1.
\]
All numerical comparisons may use $e^c\le(1-c/256)^{-256}$ for $0\le c<256$. The two finite logarithmic sums and the five initial windows are evaluated with rational arithmetic in the supplement.
\end{proof}

\paragraph{At a central Eulerian zero.}
It suffices to compare $D_{n,k}(t)$ with $A_n(t)$ at a zero of $A_{n-1}$.
Write $t=-e^{-x}$; begin with $0\le x\le(5/2)\sqrt{n-1}$, and take $n\ge128$.
Set $\rho=\sqrt{x^2+\pi^2}$, $c=x/\rho$, $s=\pi/\rho$, $z_\ell=(c+\mathrm i\ell s)^{-1}$, and $V_j=\operatorname{Re}\sum_{\ell\ge1,\ \ell\ {\rm odd}} z_\ell^j$ for $j\ge2$.
The earlier partial fraction identity gives $V_n=0$ at this zero, and the pure alphabet input $q^n$ gives $V_{n+1}$ after common rescaling.
We have $\rho<R\sqrt n$, $ns^2>3/2$, $s>6/(5\sqrt n)$, with $R=63/25$.
Bernoulli's inequality gives $|z_\ell|^{n/4}\le[1+3(\ell^2-1)/16]^{-1}$.
It follows that
\[
 \sum_{\ell\ge3,\ {\rm odd}}\ell|z_\ell|^n<1/10,\qquad
 \sum_{\ell\ge1,\ {\rm odd}}|z_\ell|^{3n/4}<11/10,\qquad
 \sum_{\ell\ge1,\ {\rm odd}}|z_\ell|^{n/2}<5/4.
\]
For each estimate separate $\ell=1,3$; above 3 use $1+3(\ell^2-1)/16\ge3\ell^2/16$ and bound the sum of $\ell^{-b}$ by $5^{-b}+5^{1-b}/[2(b-1)]$.
Since $V_n=0$, $|\operatorname{Re}z_1^n|<1/30$ and $|\operatorname{Im}z_1^n|>99/100$.
The term with $\ell=1$ in $V_{n+1}-cV_n$ therefore has magnitude greater than $(99/100)s$, while the rest costs at most $s/10$ by $|z_\ell-c|\le\ell s$. In particular $|V_{n+1}|>1/\sqrt n$.

For $d=n-h\ge2$ it remains to bound
\[
 \frac{h!}{n!}(\rho/v)^d\sum_{\ell\ge1,\ {\rm odd}}|z_\ell|^{h+1},
 \qquad v=4/5 .
\]
In the successive ranges $d=2$, $3\le d\le6$, $7\le d\le n/2$, and $d>n/2$, upper bounds are respectively
\[
 \frac{11(R/v)^2}{10(n-1)},\qquad
 \frac54\left(\frac{R\sqrt n}{v(n-5)}\right)^3,\qquad
 \frac54\left(\frac{2R}{v\sqrt n}\right)^7,\qquad
 n(4/5)^{n/2}.
\]
For the last bound use $h\ge1$, $\sum |z_\ell|^2\le\rho^2/8<n$, and $(n!/h!)^{1/d}\ge n/e$. Each bound times $\sqrt n$ is below one at $128$ and decreases for larger $n$. Multiplying by $|a_h|v^d/p_n$ and summing, Lemma~\ref{invstrong} makes the total perturbation smaller than the pure term. This proves $D_{n,k}(t)A_n(t)>0$. The identical pole representation at $x=0$ follows by continuity.

\paragraph{At an exterior Eulerian zero.}
Suppose now $x\ge(5/2)\sqrt{p}$, $p=n-1$. For $q>0$ put $g(q)=f_n(q)/(p_nq^p)$ and $w_q=q^p e^{-xq}$. Applying Lemma~\ref{invstrong} to $g(q)-q$ and its derivative gives
\[
 |g'(q)-1|\le\max_{d\ge2}(d-1)(5/(4q))^d .
\]
The upper bound is at most $125/256$ for $q\ge2$, and $25/144$ for $q\ge3$, by consecutive ratios. Also $g(1)=0$ and $39/32<g(2)<89/32$. Hence every integer increment lies in $[1/2,3]$ and every increment starting at $q\ge3$ lies in $[119/144,169/144]$.

Choose a mode $m$ of $w_q$, $q\ge1$, and let $W=w_m$. For $m=1$ the bound below is immediate and $b=2$ is the only neighbor. For $m>1$, comparison with $m-1$ gives $m\le1+(2/5)\sqrt p\le\sqrt p/2$ when $p\ge127$. Of the existing immediate neighbors let $b$ have the greater mass. The smaller neighbor has mass at most $W/7$ because the product of the two relative masses is $(1-1/m^2)^p\le e^{-4}$. At $m\le3$ it costs at most $W/50$, since $e^{-p/(2m^2)}<1/50$ for $2\le m\le3$; an absent neighbor contributes zero. At each distance $j\ge2$ the mass on either side is at most $50^{1-j}W$. Indeed the first further outward ratio is at most $e^{-p/(m-1)^2}$ on the left and $e^{-p/(m+1)^2}$ on the right by log-concavity and comparison with the mode. On the right if $m+1>\sqrt p/2$, use $\exp(p/(m+1)-x)\le e^{-\sqrt p/2}$ instead. All subsequent outward ratios only decrease.

At the Eulerian zero $\sum_q(-1)^qw_q=0$. Subtract $g(b)$ from each $g(q)$ in the corresponding signed moment. Multiply the signed moment by $\varepsilon=(-1)^m\operatorname{sgn}(m-b)$. The term at $m$ is then positive and the term at $b$ vanishes. The mass left in the three central positions, after charging the smaller neighbor, is at least $W(1/2-6/50)$ if $m\le3$, or $W(119/144-2(169/144)/7)$ if $m\ge4$. The distant terms cost at most
\[
 6W\sum_{j\ge2}(j+1)50^{1-j}=(888/2401)W,
\]
strictly smaller than both lower bounds. The same orientation works for $g(q)=q$. By \eqref{letter}, the two signed moments share the same nonzero conversion factor to $D_{n,k}(t)/p_n$ and $A_n(t)$. Thus their product is positive. For $t<-1$, Eulerian reciprocity $A_{n-1}(t)=t^{n-2}A_{n-1}(1/t)$ makes $1/t$ another Eulerian zero. Apply the argument to $(-1)^nf_n(-q)$ and $1/t$. The two evaluations acquire the same factor $t^{n-1}$, whose square is positive.

\paragraph{Completing the brackets.}
At an Eulerian zero $\xi$ we have $A_n(\xi)=\xi(1-\xi)A'_{n-1}(\xi)$. The comparison with the next Eulerian row thus yields
\[
 \frac{D_{n,k}(\xi)}{\xi}A'_{n-1}(\xi)
 =\frac{D_{n,k}(\xi)A_n(\xi)}{\xi^2(1-\xi)}>0 .
\] These strict signs alternate and are common to all the allowed cutoffs. Intermediate values give a zero in each of the $n-3$ gaps, and degree $n-3$ makes each simple and the list complete.

For $6\le n\le127$, reconstruct the same integer rows by \eqref{letter}. Isolate all $n-2$ Eulerian zeros in disjoint dyadic brackets $(l_i,r_i)$ in increasing order, indexed from $i=0$. The certificate gives signs $-s_i,s_i$ for $A_{n-1}$ at their endpoints, where $s_i=(-1)^{n-3-i}$. For every quotient the sign is $s_i$ at \emph{both} endpoints. Every gap $(r_i,l_{i+1})$ therefore contains a zero of every quotient. This verifies 7,869 pairs through 3,904 cutoff representatives; all $k>n/2$ give the single-cycle row. The argument by degree just given also applies to the finite certificates. Since the sign inequalities are linear in the quotient and a nonzero nonnegative combination has positive leading coefficient, the conclusion extends to the mixtures in Theorem~\ref{gapthm}.
